\chapter{Estado del arte}\label{cap:estado}

Este capítulo revisa tres líneas de trabajo que la tesis cruza: la detección de lavado de
dinero con grafos, en particular sobre Elliptic (Sección \ref{sec:ea_aml}); la explicabilidad
aplicada a ese problema (Sección \ref{sec:ea_xai_aml}); y la evaluación de la estabilidad de
las explicaciones, en especial bajo desbalance de clases (Sección \ref{sec:ea_estab}). Al
final se precisa el vacío que aborda el trabajo (Sección \ref{sec:ea_vacio}).

\section{Detección de lavado de dinero con grafos}\label{sec:ea_aml}

Las revisiones sobre aprendizaje automático para la prevención del lavado de dinero identifican técnicas de análisis de vínculos, modelado del comportamiento, puntuación de riesgo y detección de anomalías \cite{Chen2018AMLReview}. Kute et al. observan que el aprendizaje profundo se ha usado poco en este dominio, sobre todo por la falta de interpretabilidad de sus decisiones, y que faltan estudios que lo combinen con explicabilidad \cite{Kute2021AMLReview}. La detección de fraude financiero con GNN ya cuenta con una revisión sistemática dedicada \cite{Motie2024FraudGNN}.

Elliptic se publicó junto con los primeros experimentos de GCN para este problema
\cite{Weber2019Elliptic}. Sus autores compararon regresión logística, bosques aleatorios y GCN,
y observaron que el desempeño en la prueba se desplomaba a partir del paso 43, tras el cierre
de un mercado ilegal. Ese hallazgo es la base de la partición que usa esta tesis. Desde
entonces el conjunto de datos se ha convertido en la referencia para la detección de
transacciones ilícitas en Bitcoin. En criptomonedas se han propuesto modelos de \textit{gradient boosting} que superan a los bosques aleatorios y cuya adaptación reduce el efecto del cambio de concepto \cite{Vassallo2021Boosting}. Sobre Elliptic se han propuesto modelos que combinan GCN con memorias de largo y corto
plazo para aprovechar el orden temporal, junto con esquemas de aprendizaje activo
\cite{Alarab2023Graph}; representaciones de nodos aprendidas de forma autosupervisada
\cite{Lo2023Inspection}; y arquitecturas temporales basadas en transformadores sobre grafos
\cite{Lin2026Detecting}. Otras propuestas modelan el grafo de transacciones de forma dinámica: Wan et al. construyen instantáneas del grafo en el tiempo y las combinan con memorias de largo y corto plazo, centrándose en la etapa de estratificación \cite{Wan2024Memory}, y Wang et al. proponen GNN dinámicas semisupervisadas para fraude en cadenas de bloques \cite{Wang2025CoSemiGNN}. Mohan et al. modelan la detección en Bitcoin como clasificación de
nodos en redes dinámicas, combinando convoluciones de grafos que evolucionan en el tiempo con
bosques de decisión neuronales \cite{DTA2022Evolving}, y Asiri et al. comparan regresión
logística, máquinas de soporte vectorial, bosques aleatorios, redes LSTM y GCN sobre Elliptic
\cite{SciRep2025GCN}.

El cierre del mercado en el paso 43 es un problema reconocido. Alarab y Prakoonwit señalan que
todos los algoritmos de aprendizaje han fallado en capturar ese evento con las variables
originales de Elliptic, y proponen un modelo recurrente de grafos que hace predicciones
secuenciales apoyándose en las transacciones etiquetadas más recientes
\cite{Alarab2023Recurrent}. Para ello añaden a cada transacción dos variables, el número de
vecinas antecedentes lícitas y el de ilícitas, y con ellas reportan un F1 superior a 0,80 en el
paso 43. Ese enfoque supone que se dispone de etiquetas recientes durante la
operación, un supuesto distinto del de esta tesis, que separa estrictamente el entrenamiento de
la evaluación. En el ámbito forense,
Pocher et al. analizan cómo los métodos de aprendizaje automático y de análisis del grafo de
transacciones se aplican a la prevención del lavado en criptomonedas y qué implicaciones
tiene su uso por autoridades y empresas \cite{Pocher2023Anomalous}.

Dos observaciones de esta literatura son relevantes para la tesis. La primera es que la
mayoría de los trabajos reporta el desempeño con la partición temporal habitual o con
particiones aleatorias, y pocos discuten el efecto del cierre del mercado sobre la prueba. La
segunda es que el desbalance de clases se trata como un obstáculo para el rendimiento. Lawal
et al. examinan explícitamente las limitaciones de las GNN sobre Elliptic bajo ese
desbalance. En su experimento, que trata las transacciones sin etiqueta como una tercera clase,
el modelo alcanzó una exactitud del 100~\% en entrenamiento y validación, pero en la prueba,
de ocho nodos, predijo la clase mayoritaria para todos \cite{Lawal2025Dataset}. Su resumen
menciona pesos por clase, \textit{focal loss} y GNNExplainer, pero el texto no presenta
resultados de esas variantes, por lo que aquí se cita solo como ejemplo de una exactitud
engañosa. Ninguno de estos trabajos mide qué ocurre con las
explicaciones cuando cambia el desbalance.

\section{Explicabilidad en la detección de lavado de dinero}\label{sec:ea_xai_aml}

La falta de interpretabilidad es, según Kute et al., la razón principal por la que el aprendizaje profundo se ha usado poco en la detección de lavado \cite{Kute2021AMLReview}. Una revisión reciente reúne los
trabajos sobre detección de lavado en criptomonedas que incorporan explicabilidad
\cite{Unagar2025Survey}, y se han propuesto marcos de GNN explicables para detectar transacciones financieras ilícitas, como el de He et al., construido sobre TAGCN para integrar información de vecindarios de orden superior \cite{He2026Detection}. En la mayoría de estos trabajos la
explicación se presenta como un resultado cualitativo, por ejemplo las variables o los vecinos
más importantes para algunos casos, y no se evalúa si la explicación se mantiene al repetirla.

En el campo general de la explicabilidad de GNN, las taxonomías clasifican los métodos según lo que explican (variables, aristas o subgrafos) y cómo lo obtienen (gradientes, perturbaciones, sustitutos o descomposición) \cite{Yuan2023Taxonomic}. Las
evaluaciones comparativas han mostrado que los explicadores discrepan entre sí. Agarwal et al.
construyeron un generador de grafos sintéticos con explicaciones de referencia, capaz de
variar el tamaño, la distribución de grados y la homofilia del grafo, y una biblioteca para
evaluar explicadores sobre esos grafos y sobre grafos reales \cite{Agarwal2023Evaluating}.
Entre sus medidas incluyen la estabilidad, entendida como el parecido entre la explicación de un
grafo y la de una versión con perturbaciones mínimas, y encontraron que la estabilidad de las
explicaciones sobre variables se parece más a la de explicaciones aleatorias que la de las
explicaciones sobre nodos. Longa et al. probaron 12 explicadores
sobre 8 arquitecturas de paso de mensajes y 6 conjuntos de datos, midiendo la plausibilidad y la
fidelidad de las explicaciones, y analizaron por qué unas
arquitecturas resultan más explicables que otras y qué explicador conviene en cada caso
\cite{Longa2024Explaining}.

\section{Estabilidad de las explicaciones y desbalance de clases}\label{sec:ea_estab}

Las revisiones sobre la evaluación cuantitativa de la explicabilidad incluyen la estabilidad,
o consistencia, entre las propiedades que una explicación debe cumplir, junto a la fidelidad y
la plausibilidad \cite{Nauta2023Anecdotal, Bodria2023Benchmarking}. Se han propuesto métricas
específicas para medirla \cite{Gawantka2024Explanations}. Adebayo et al. mostraron que algunas
explicaciones visualmente convincentes apenas cambian cuando se aleatorizan los pesos del
modelo, y propusieron esa aleatorización como prueba de que una explicación depende de lo que
el modelo aprendió \cite{Adebayo2018Sanity}.

La relación entre desbalance de clases y explicaciones se ha estudiado en datos tabulares. Rai
et al. diseñaron un marco comparativo para evaluar la confiabilidad y la consistencia de LIME,
SHAP y gráficos de dependencia parcial en la predicción del riesgo de cáncer de pulmón con
historias clínicas electrónicas \cite{Rai2025Evaluating}. Entrenaron modelos con proporciones de
casos del 40~\%, 30~\%, 20~\%, 10~\%, 5~\% y 1~\% y con la distribución original, y midieron la
consistencia como el acuerdo, con el índice de Jaccard de las 10 variables principales y un
acuerdo de rangos, entre las explicaciones de cada modelo y las de un modelo de referencia
entrenado con datos balanceados. Esa consistencia baja a medida que aumenta el desbalance. Los
autores distinguen expresamente su medida de la estabilidad entre entrenamientos repetidos que
estudian Chen et al.
Fujiwara et al. proponen combinar destilación de conocimiento y remuestreo para mejorar tanto
el rendimiento como la estabilidad de las explicaciones en clasificación desbalanceada
\cite{Fujiwara2024Stability}. Sobre Elliptic, Alarab y Prakoonwit compararon técnicas de
remuestreo con un bosque aleatorio entrenado sobre las transacciones etiquetadas, y midieron la
importancia de las variables por permutación antes y después de remuestrear. Con la prueba de
Wilcoxon encontraron cambios significativos solo con algunas técnicas y conjuntos: en la prueba
de Elliptic, con SMOTE-SF ($p = 0{,}001$); en el entrenamiento, con ninguna
\cite{Alarab2022Resampling}. Ese trabajo examina si cambian las variables importantes, no la
estabilidad de la explicación al repetirla.

El trabajo más cercano a la pregunta de esta tesis es el de Chen et al. sobre calificación
crediticia. Partiendo de que la investigación sobre desbalance se había concentrado en su
efecto sobre la exactitud y no sobre la interpretabilidad, evaluaron la estabilidad de LIME y
SHAP sobre conjuntos de datos hipotecarios con desbalance creciente, y encontraron que las
explicaciones se vuelven menos estables a medida que aumenta el desbalance, resultado que
repitieron en otros dos conjuntos de datos de crédito \cite{Chen2023CreditImbalance}. Mammadov
et al. propusieron un artefacto de diseño para evaluar de forma controlada el efecto del
desbalance y del muestreo sobre la estabilidad de SHAP y LIME, con medidas de acuerdo de rangos
y de variabilidad, sobre 1.600 simulaciones \cite{TEM2026Stability}. Estos trabajos muestran que
el desbalance y las técnicas que lo compensan pueden afectar las explicaciones, y en el caso de
Chen et al. ofrecen una hipótesis concreta, la de que más desbalance produce explicaciones
menos estables, que esta tesis contrasta en un dominio distinto. Ninguno trata datos en grafo
ni explicadores de GNN.

\section{Síntesis comparativa}

La Tabla \ref{tab:ea} resume los trabajos más cercanos según cuatro rasgos: si usan grafos, si
explican las predicciones, si miden la estabilidad de las explicaciones y si varían el
desbalance de forma controlada.

\begin{table}[ht]
\centering
\caption{Trabajos relacionados según su cobertura de los cuatro rasgos que combina la tesis,
según su texto completo, salvo Chen et al., Mammadov et al. y He et al., para los que se usó el
resumen.
\checkmark: lo hace; --: no lo hace. Desbalance controlado significa variar el nivel de
desbalance, no solo la técnica que lo compensa. $^a$ Estabilidad frente a perturbaciones
mínimas de la entrada. $^b$ Acuerdo con las explicaciones de un modelo balanceado, no entre
réplicas.}
\label{tab:ea}
\small
\begin{tabularx}{\textwidth}{Xcccc}
\toprule
Trabajo & Grafos & Explicación & Estabilidad medida & Desbalance controlado \\
\midrule
Weber et al. \cite{Weber2019Elliptic} & \checkmark & -- & -- & -- \\
Alarab et al. \cite{Alarab2023Graph, Alarab2023Recurrent} & \checkmark & \checkmark & -- & -- \\
Lawal et al. \cite{Lawal2025Dataset} & \checkmark & -- & -- & -- \\
Alarab y Prakoonwit \cite{Alarab2022Resampling} & -- & \checkmark & -- & -- \\
He et al. \cite{He2026Detection} & \checkmark & \checkmark & -- & -- \\
Agarwal et al. \cite{Agarwal2023Evaluating} & \checkmark & \checkmark & \checkmark$^a$ & -- \\
Longa et al. \cite{Longa2024Explaining} & \checkmark & \checkmark & -- & -- \\
Chen et al. \cite{Chen2023CreditImbalance} & -- & \checkmark & \checkmark & \checkmark \\
Mammadov et al. \cite{TEM2026Stability} & -- & \checkmark & \checkmark & \checkmark \\
Rai et al. \cite{Rai2025Evaluating} & -- & \checkmark & \checkmark$^b$ & \checkmark \\
Esta tesis & \checkmark & \checkmark & \checkmark & \checkmark \\
\bottomrule
\end{tabularx}
\end{table}

\section{Vacío de investigación}\label{sec:ea_vacio}

De la revisión se desprende que (1) los trabajos sobre Elliptic se concentran en el
rendimiento y tratan la explicación de forma cualitativa; (2) las evaluaciones sistemáticas de
explicadores de GNN usan sobre todo grafos sintéticos o de otros dominios, miden la estabilidad,
cuando la miden, frente a perturbaciones de la entrada y no entre réplicas, y no varían el
desbalance; y (3) los estudios sobre desbalance y explicaciones usan datos tabulares. No se
encontró un estudio que mida la estabilidad de varios explicadores de GNN sobre un problema
real de lavado de dinero mientras se controlan, en un diseño factorial con réplicas
independientes, el nivel de desbalance, la arquitectura y la técnica de compensación. Tampoco
se encontró uno que verifique, con un control de pesos aleatorios, si la estabilidad medida en
ese problema refleja lo que el modelo aprendió. Esta tesis aborda ambas carencias.
